\section{The cost of a prompt token}
\label{sec:cost}

\paragraph{Recompute, restore, hit.} To price the three ways a prompt token can be served, we sent eight prompts of
the same length at once with one output token each, three times, after flushing the GPU cache, and subtracted idle
power (Table~\ref{tab:cost}). Recomputing costs 56--73\,mJ per prompt token for Qwen3-32B in BF16, rising with length
by 1.26$\times$ from about 5k to 20k tokens; Eq.~\eqref{eq:recompute} predicts 1.24$\times$. Restoring the same
prompts from SGLang's 160\,GB host tier costs 0.27--0.68\,mJ per token, at or below the 0.5--1.3\,mJ predicted
by Eq.~\eqref{eq:restore}, with 99.98--100\% of tokens served from the tier. A GPU hit costs 0--0.3\,mJ. Recomputing
thus costs about 110, 160 and 220 times as much as restoring at the three lengths; the gap widens
2$\times$ with length, more than the 1.24$\times$ of $e_r$ alone, because the measured restore cost per token also falls with length. The median time to first token for eight 20k-token prompts was 1.36\,s restored and 15.5\,s
recomputed.
As a control, the same procedure with no host tier left 0.1--0.4\% of tokens cached and energy within 4\% of
recompute.

These figures are above-idle GPU board energy for a whole batch, so they include the one output token's decode and
scheduling: they bound the GPU-side cost of reading KV back from the host tier from above; the write to the host during the first send is not separable from recompute and is not included. Host CPU and memory energy are not included; Mao and
Zhou~\cite{mao2026offloading} bounded the host-side difference of offloading on a similar H200 host at under 1\% of the
GPU saving. Mao and Zhou also estimated that moving a KV byte is at least 400$\times$ cheaper than recomputing it; the
llm-d research agenda~\cite{sisodia2026toward} asks when transfer is greener than recompute. Our numbers answer it
directly for host-memory restore over PCIe, on one card and one model.

\begin{table}[t]
\centering\small
\caption{Above-idle GPU energy per prompt token, Qwen3-32B (BF16) on an H200, three repetitions of eight prompts. Ranges span the three repetitions (recompute: those of both the host-tier and the control runs). Hit and restore phases last 0.1--0.2\,s and
0.3--1.4\,s, respectively, near the meter's 0.1\,s sampling, so their small values are approximate (two hit readings are slightly
negative); the ratios are of the means (85--260$\times$ per repetition).}
\label{tab:cost}
\begin{tabular}{lrrr}
\toprule
Prompt length (tokens) & $\sim$5k & $\sim$10k & $\sim$20k \\
\midrule
Recompute (mJ) & 56--58 & 62--64 & 73 \\
Restore from host tier (mJ) & 0.42--0.68 & 0.27--0.46 & 0.28--0.37 \\
GPU hit (mJ) & 0--0.3 & 0--0.3 & 0--0.3 \\
Recompute / restore & $\sim$110$\times$ & $\sim$160$\times$ & $\sim$220$\times$ \\
\bottomrule
\end{tabular}
\end{table}

\paragraph{Holding is free.} Keeping KV in GPU memory costs no measurable power. Idle power with the cache filled to
capacity (261k prompt tokens sent against 245--251k of capacity; the first fill prompt was still cached afterwards on
SGLang and had been evicted on vLLM) differed from idle power with an empty cache by
$-0.08\%$ to $-0.15\%$ on vLLM and $-0.13\%$ to $-0.16\%$ on SGLang (two readings each), after equal temperature settling. The energy question is therefore mostly
where to keep KV: holding is free, rebuilding costs 56--73\,mJ per token, and moving it costs well under
1\,mJ of GPU energy.

\paragraph{The penalty grows with load.} With one user, a cache miss on a 16k-token prompt cost 1.46--1.48$\times$ the
energy per output token of a hit; at 64 concurrent users it cost 10.1--10.6$\times$, on both engines. Misses cost most
when the server is busy; a likely cause, which we did not measure, is that recomputed prefill displaces batched
decoding. The 64-user figure is an upper bound on recompute alone: there the hit requests share
one 16k-token document, while the miss requests need 64 separate contexts, about four times the GPU's KV capacity, so
the miss case also loses batch size to memory and queueing.
